\section{Related Work and Positioning}

\subsection{Graph-based fraud detection}
Graph fraud research has largely advanced by changing the representation or learning rule. GraphConsis regularizes inconsistent neighborhoods \cite{liu2020graphconsis}; CARE-GNN addresses camouflaged fraudsters through relation-aware aggregation \cite{dou2020caregnn}; and PC-GNN couples class-aware sampling with neighborhood aggregation for imbalance \cite{liu2021pcgnn}. BWGNN targets anomalous graph signals in the spectral domain \cite{tang2022bwgnn}. Directed-multigraph models preserve edge direction and multiplicity that homogeneous projections can erase \cite{egressy2024directed,wu2023grande}, while FraudGT develops an efficient transformer for financial graphs \cite{lin2024fraudgt}. These methods sit on the foundations provided by GCN, GraphSAGE, graph attention, and graph isomorphism networks \cite{kipf2017gcn,hamilton2017graphsage,velickovic2018gat,xu2019gin}.

This modeling work explains why graph construction cannot be treated as neutral. A fraud detector may consume account links, transaction nodes, directed transaction edges, relation types, or aggregates computed from a historical subgraph. Two recent results are especially relevant. GADBench finds that simple non-GNN methods can be competitive across supervised graph-anomaly tasks \cite{tang2023gadbench}; the Graph Feature Preprocessor obtains efficient financial-crime models by composing graph-derived features with tabular learners \cite{blanusa2024gfp}. We therefore use logistic regression and histogram gradient boosting as substantive baselines, and treat edge attributes, neighborhood summaries, degree caps, and recency windows as testable constructions. FraudShiftBench does not propose another universal detector. It asks which detector conclusion is supported under a stated information contract.

\subsection{Financial and graph-anomaly benchmarks}
Elliptic, DGraphFin, and IBM AML-Data cover markedly different tasks: illicit Bitcoin transaction nodes, anomalous account nodes, and laundering transaction edges \cite{weber2019aml,huang2022dgraph,altman2023ibmaml}. Reviews of graph-based fraud detection already document the diversity of entity definitions, relations, and operational assumptions \cite{pourhabibi2020survey}. Broader anomaly benchmarks expand method and dataset coverage \cite{tang2023gadbench}; BAG adds dynamic graph tasks \cite{hua2026bag}; and GAD in the Wild studies scale, rarity, and missing attributes under deployment-oriented settings \cite{zhou2026wild}. These works preclude a credible claim that ours is the first dynamic or realistic graph-anomaly benchmark.

The remaining gap is narrower. Existing benchmarks do not jointly type the temporal unit, graph visibility, construction, selection rule, investigation budget, and compute envelope of a fraud claim, then check that the necessary predictions and paired cells exist. Dataset breadth alone cannot resolve that gap. FraudShiftBench consequently keeps the three task units separate and reports conditional comparisons within each compatible scope.

\subsection{Temporal evaluation and distribution shift}
OGB established common graph tasks, evaluators, and dataset splits at scale \cite{hu2020ogb}; GOOD isolates graph distribution shifts \cite{gui2022good}. For dynamic link prediction, Poursafaei et al. show that evaluation choices such as negative sampling and inductive status materially affect conclusions \cite{poursafaei2022better}. TGB, TGB~2.0, and BenchTemp extend standardized evaluation across temporal, heterogeneous, and knowledge graphs, with attention to efficiency and temporal order \cite{huang2023tgb,gastinger2024tgb2,huang2024benchtemp}.

Fraud detection adds an unusual separation between structural and label availability. An account or transaction may be visible when scored even though its eventual fraud label is not. The historical graph used to construct node features may also cover a different interval from the labels used to fit a classifier. Our deployment contract records those choices separately. It complements general temporal benchmarks by asking a fraud-specific question: which covariates, endpoints, edges, labels, and decision capacities are available at each stage?

\subsection{Validity, documentation, and benchmark governance}
Standard splits can conceal sampling assumptions \cite{gorman2019splits}, while feature leakage and repeated model selection can invalidate an otherwise reproducible pipeline \cite{kaufman2012leakage,cawley2010selection,kapoor2023leakage}. The benchmark lottery describes a related external-validity problem: performance depends on which benchmark happens to stand in for the intended task \cite{dehghani2021lottery}. BetterBench evaluates benchmark quality systematically rather than accepting a leaderboard at face value \cite{reuel2024betterbench}; recent work also argues that broad state-of-the-art claims require commensurately broad evidence \cite{oh2026sota}.

Documentation frameworks solve another part of the problem. Data Statements, Datasheets, Data Cards, and Model Cards make intended use, composition, and evaluation conditions visible \cite{bender2018data,gebru2021datasheets,pushkarna2022datacards,mitchell2019modelcards}. BenchmarkCards and Eval Factsheets extend structured reporting to benchmarks and evaluations \cite{sokol2025benchmarkcards,bordes2025evalfacts}. Measurement work cautions that even well-documented metrics can fail to represent the construct of interest \cite{jacobs2021measurement}. FraudShiftBench builds on these practices but changes the unit of validation: a typed empirical claim must be supported by complete evidence under its requested deployment scope. The validator does not certify scientific truth; it checks the declared relation among scope, predictions, pairing, integrity, and feasibility.

\subsection{Ranking, calibration, and constrained decisions}
Rare-event screening requires more than a single discrimination score. Precision--recall curves are informative under severe imbalance \cite{saito2015pr}, but AUPRC still describes score ordering rather than the alert policy that investigators execute. Probability calibration can improve score interpretation \cite{guo2017calibration}, and prior-shift correction can adapt posterior probabilities when class prevalence changes \cite{saerens2002prior}. Selective classification studies risk as a function of accepted coverage \cite{geifman2017selective}; cost-sensitive learning attaches explicit loss to different errors \cite{bahnsen2015cost}; and feedback-guided anomaly discovery models inspection of a ranked queue \cite{siddiqui2018feedback}.

This motivates a strict distinction between \emph{rank claims} and \emph{decision claims}. AUROC and AUPRC depend on ordering. F1, precision or recall at a review budget, and cost-sensitive risk depend on a threshold or capacity. A strictly increasing recalibration preserves the former while it may change the latter. FraudShiftBench records the decision rule in the contract so that an operational gain cannot be described as a repaired ranking.

\begin{table*}[t]
\caption{Closest benchmark families and the distinctions needed here. ``No'' means that the feature is not a stated benchmark axis; it does not question the rigor of the cited work.}
\label{tab:related-comparison}
\centering
\footnotesize
\setlength{\tabcolsep}{3.2pt}
\begin{tabularx}{\textwidth}{@{}p{0.16\textwidth}c p{0.19\textwidth}p{0.10\textwidth}p{0.11\textwidth}p{0.10\textwidth}X@{}}
\toprule
Work & Temporal & Graph contract & Capacity & Resources & Predictions & Claim support \\
\midrule
GADBench & No & Limited & No & No & Benchmark & No \\
TGB / BenchTemp & Yes & Inductive tasks & No & Efficiency & Benchmark & No \\
BAG & Yes & Task/model breadth & No & Scale & Benchmark & No \\
GAD in the Wild & Partial & Scale/missing attrs. & No & Scale & Limited & No \\
BetterBench / Eval Factsheets & Generic & Documented & Documented & Documented & Metadata & Descriptive \\
FraudShiftBench & Yes & Visibility + construction & Top-$K$ + cost & Typed status & Per-seed & Executable \\
\bottomrule
\end{tabularx}
\end{table*}

